\documentclass[11pt]{article}

\usepackage[T1]{fontenc}
\usepackage{amsmath,amssymb,mathtools}
\usepackage[margin=1in]{geometry}
\usepackage{microtype}

\newcommand{\R}{\mathbb R}
\newcommand{\cA}{\mathcal A}
\newcommand{\cG}{\mathcal G}
\newcommand{\cO}{\mathcal O}
\newcommand{\cQ}{\mathcal Q}
\newcommand{\cV}{\mathcal V}
\newcommand{\ev}{\operatorname{ev}}
\newcommand{\supp}{\operatorname{supp}}
\newcommand{\weakLthree}{L^{3,\infty}}
\newcommand{\LorentzDual}{L^{3/2,1}}

\title{Adjoint-pressure cost of one Besov regeneration bridge}
\author{}
\date{23 July 2026}

\begin{document}
\maketitle

\paragraph{Scope.}
This note records a conditional a priori estimate on the smooth physical
genealogy converging to the ancient weak-\(L^3\) profile.  It does not assume
that a rough hull state has a uniquely defined endpoint Oseen adjoint, and it
does not establish a finite sum of the costs over all regeneration bridges.

\section*{1. Smooth finite-horizon adjoint estimate}

Here \(\nu\) denotes the viscosity of the trajectory under discussion.
For the normalised outer profile in Section 2, \(\nu=1\).

Let \(u\) be a smooth, finite-energy, divergence-free Navier--Stokes
solution on
\(\R^3\times[-H,0]\), with viscosity \(\nu>0\).  Let \(M>0\) and assume
\begin{equation}
  \sup_{-H\le t\le0}
  \|u(t)\|_{\weakLthree(\R^3)}
  \le M.
  \label{eq:weak-L3}
\end{equation}
Put \(b(\tau)=u(-\tau)\) and \(g=u(0)\).  Given
\(\psi\in C^\infty_{c,\sigma}(\R^3;\R^3)\), let
\((a,\pi^*)\) solve the solenoidal Oseen adjoint
\begin{equation}
  \boxed{
  \begin{aligned}
    \partial_\tau a-\nu\Delta a-b\cdot\nabla a+\nabla\pi^*&=0,
      &&0<\tau<H,\\
    \nabla\cdot a&=0,\\
    a(0)&=\psi .
  \end{aligned}}
  \label{eq:adjoint}
\end{equation}
Equivalently,
\[
  \partial_\tau a-\nu\Delta a-\mathbb P(b\cdot\nabla a)=0,
\]
where \(\mathbb P\) is the Leray projection.  In the reversed time
variable, the primal velocity satisfies
\[
  \partial_\tau b=-\nu\Delta b+\mathbb P(b\cdot\nabla b).
\]
The drift terms cancel after integration by parts, and therefore
\begin{equation}
  \boxed{\langle b(T),a(T)\rangle=\langle g,\psi\rangle}
  \qquad(0\le T\le H).
  \label{eq:pairing}
\end{equation}
The adjoint energy identity is
\begin{equation}
  \boxed{
  \|a(T)\|_2^2
+2\nu\int_0^T\|\nabla a(\tau)\|_2^2\,d\tau
=\|\psi\|_2^2.}
  \label{eq:energy}
\end{equation}

Set
\[
  B(T):=\sup_{0\le\tau\le T}\|a(\tau)\|_1.
\]
If \(B(T)<\infty\), Nash's inequality and \eqref{eq:energy} give
\begin{equation}
  \|a(T)\|_2
  \le C_{\mathrm N}B(T)(\nu T)^{-3/4}.
  \label{eq:nash}
\end{equation}
Real interpolation between \(L^1\) and \(L^2\), with interpolation
parameter \(2/3\), then gives
\begin{align}
  \|a(T)\|_{\LorentzDual}
  &\le C_{\mathrm{int}}
     \|a(T)\|_1^{1/3}\|a(T)\|_2^{2/3} \notag\\
  &\le C_*B(T)(\nu T)^{-1/2}.
  \label{eq:interpolation}
\end{align}
On the other hand, Lorentz duality and \eqref{eq:pairing} imply
\[
  |\langle g,\psi\rangle|
  \le C_{\mathrm{Lor}}M
      \|a(T)\|_{\LorentzDual}.
\]
Consequently, with
\(C_{\mathrm{adj}}:=C_{\mathrm{Lor}}C_*\),
\begin{equation}
  \boxed{
  B(T)\ge
  \frac{|\langle g,\psi\rangle|}
       {C_{\mathrm{adj}}M}\sqrt{\nu T}.}
  \label{eq:L1-growth}
\end{equation}
The assertion remains valid in the extended sense if \(B(T)=\infty\).

The vector Kato inequality applied to \eqref{eq:adjoint} gives
\[
  \frac{d}{d\tau}\|a(\tau)\|_1
  \le \|\nabla\pi^*(\tau)\|_1.
\]
If the pressure-history integral is finite, this follows rigorously by
testing a smooth approximation of \(|a|\) against a spatial cutoff,
letting the cutoff radius tend to infinity, and then removing the modulus
regularisation; the smooth finite-energy drift is bounded on every finite
horizon.  If that integral is infinite, the estimate below is automatic.
Diffusion is \(L^1\)-contractive and the divergence-free drift integrates
to zero.  Combining this inequality with \eqref{eq:L1-growth} proves the
finite-horizon estimate
\begin{equation}
  \boxed{
  \int_0^T
  \|\nabla\pi^*(\tau)\|_{L^1(\R^3)}\,d\tau
  \ge
  \frac{|\langle g,\psi\rangle|}
       {C_{\mathrm{adj}}M}\sqrt{\nu T}
  -\|\psi\|_{L^1(\R^3)}.}
  \label{eq:finite-horizon}
\end{equation}

\section*{2. Lower-relaxed cost on a physical genealogy}

Let \(X\) be an ancient trajectory state with terminal trace
\(g=\ev_0X\).  An admissible smooth physical genealogy for \(X\) is a
sequence
\[
  \cG=\{(u_n,H_n)\}_{n\ge1}
\]
such that:
\begin{enumerate}
  \item \(u_n\) is a smooth finite-energy Navier--Stokes solution on
        \(\R^3\times[-H_n,0]\);
  \item \(H_n\to\infty\);
  \item
        \[
          \sup_n\sup_{-H_n\le t\le0}
          \|u_n(t)\|_{\weakLthree}\le M;
        \]
  \item \(u_n\to X\) strongly in local spacetime \(L^3\) on every
        compact negative-time cylinder, with the corresponding local weak
        pressure convergence;
  \item \(u_n(0)\to g\) in \(\mathcal D'(\R^3)\).
\end{enumerate}
For the outer profile, such a sequence is obtained from its physical
spacetime diagonal \(V_n\) at radii \(\rho_n\to0\).  Choose a countable
determining family
\(\{\chi_k\}\subset C_c^\infty(\R^3;\R^3)\), dense in
\(L^{3/2,1}\) on a compact exhaustion.  Terminal weak continuity permits
the diagonal choice \(0<\varepsilon_n<1/n\) such that
\[
  \left|
  \langle V_n(-\varepsilon_n)-V_n(0),\chi_k\rangle
  \right|<\frac1n
  \qquad(k\le n).
\]
The uniform weak-\(L^3\) ceiling and Lorentz duality extend this
convergence to every compact smooth test, hence
\(V_n(-\varepsilon_n)-V_n(0)\to0\) in \(\mathcal D'\).  Set
\[
  u_n(s):=V_n(s-\varepsilon_n),
  \qquad
  H_n:=\frac{\nu_{\mathrm{phys}}T^*}{2\rho_n^2}.
\]
The translated interval lies strictly before the first singular time,
so \(u_n\) is smooth; \(H_n\to\infty\), critical scaling preserves the
weak-\(L^3\) ceiling, and the vanishing translation preserves the local
trajectory limit.  Here \(\nu_{\mathrm{phys}}\) is the original physical
viscosity; the profiles \(u_n\) have unit viscosity.

For \(T\le H_n\), denote by \(\pi^*_{n,\psi}\) the pressure in
\eqref{eq:adjoint} driven by \(b_n(\tau)=u_n(-\tau)\).  Define
\begin{equation}
  \boxed{
  \mathfrak q^\cG_\psi(X)
  :=
  \liminf_{T\to\infty}
  \ \liminf_{\substack{n\to\infty\\H_n\ge T}}
  \frac{1}{\sqrt{\nu T}}
  \int_0^T
  \|\nabla\pi^*_{n,\psi}(\tau)\|_{L^1(\R^3)}
  \,d\tau .}
  \label{eq:cost}
\end{equation}
Taking first \(n\to\infty\) and then \(T\to\infty\) in
\eqref{eq:finite-horizon} yields
\begin{equation}
  \boxed{
  \mathfrak q^\cG_\psi(X)
  \ge
  \frac{|\langle\ev_0X,\psi\rangle|}
       {C_{\mathrm{adj}}M}.}
  \label{eq:cost-lower}
\end{equation}
The lower bound is independent of the admissible genealogy.  The value of
\(\mathfrak q^\cG_\psi(X)\) is not asserted here to be intrinsic to an
arbitrary rough hull state.

For reference, if
\[
  \psi_\lambda(x):=\lambda^{-2}\psi(x/\lambda),
\]
then
\[
  \|\psi_\lambda\|_1=\lambda\|\psi\|_1,\qquad
  \|\psi_\lambda\|_2^2=\lambda^{-1}\|\psi\|_2^2,\qquad
  \|\psi_\lambda\|_{\LorentzDual}
  =\|\psi\|_{\LorentzDual}.
\]
Writing \(\mathcal S_\lambda\) for the Navier--Stokes dilation, the
corresponding scaled genealogy satisfies
\begin{equation}
  \boxed{
  \mathfrak q^\cG_{\psi_\lambda}(X)
  =
  \mathfrak q^{\mathcal S_\lambda\cG}_\psi
  (\mathcal S_\lambda X).}
  \label{eq:scaling}
\end{equation}

\section*{3. Cost of one Besov regeneration bridge}

Let \(X_*\) be the ancient profile, let
\(\Phi_\theta=\mathcal S_{e^\theta}\), and choose
\(\varphi\in C^\infty_{c,\sigma}(\R^3;\R^3)\).  Put
\begin{equation}
  z(\theta)
  :=
  \left\langle
    \ev_0(\Phi_\theta X_*),\varphi
  \right\rangle ,
  \qquad
  \cA\varphi:=2\varphi+x\cdot\nabla\varphi .
  \label{eq:event-observable}
\end{equation}
The exact dilation derivative is
\begin{equation}
  z'(\theta)
  =
  -\left\langle
    \ev_0(\Phi_\theta X_*),\cA\varphi
  \right\rangle .
  \label{eq:dilation-derivative}
\end{equation}
Suppose that consecutive event times satisfy
\[
  z(\theta_m)\ge c_0,\qquad
  z(\theta_{m+1})\ge c_0,\qquad
  \theta_{m+1}-\theta_m\ge1,
\]
and put \(I_m=[\theta_m,\theta_{m+1}]\).  Define
\[
  \cV_m:=\int_{I_m}|z'(\theta)|\,d\theta,
  \qquad
  \cO_m:=
  \bigl|\{\theta\in I_m:z(\theta)\ge c_0/2\}\bigr|.
\]
Let
\(\cG_\theta:=\mathcal S_{e^\theta}\cG\).
To avoid assuming measurability of the genealogy-relaxed functional,
write \(\int^{\!*}\) for the lower Lebesgue integral and define
\begin{align}
  \cQ_m
  :={}&
  \int_{I_m}^{\!*}
  \mathfrak q^{\cG_\theta}_{\cA\varphi}
      (\Phi_\theta X_*)\,d\theta \notag\\
  &+
  \int_{I_m}^{\!*}
  \mathbf 1_{\{z(\theta)\ge c_0/2\}}
  \mathfrak q^{\cG_\theta}_{\varphi}
      (\Phi_\theta X_*)\,d\theta .
  \label{eq:bridge-cost}
\end{align}
Equations \eqref{eq:cost-lower} and
\eqref{eq:dilation-derivative} imply
\begin{equation}
  \boxed{
  \cQ_m
  \ge
  \frac{\cV_m+(c_0/2)\cO_m}
       {C_{\mathrm{adj}}M}
  \ge
  \frac{c_0}{2C_{\mathrm{adj}}M}>0.}
  \label{eq:single-event}
\end{equation}
Indeed, either \(z\ge c_0/2\) throughout \(I_m\), in which case
\(\cO_m\ge1\), or \(z\) drops below \(c_0/2\) and returns to at least
\(c_0\), in which case \(\cV_m\ge c_0\).

\paragraph{What remains open.}
Estimate \eqref{eq:single-event} is a lower bound for one regeneration
bridge.  The adjoint histories associated with different scales overlap,
and no estimate proved here bounds
\(\sum_{m=1}^N\cQ_m\) uniformly in \(N\).  Such an event-index summation,
or a uniform \(o(\sqrt T)\) upper bound contradicting
\eqref{eq:finite-horizon}, would be the new rigidity input needed to close
this conditional branch.

\end{document}
